\scanpage{9}
Lie, donc $\unsure{G|C}$ opère dans $\unsure{\mathfrak{k}|k}$. Soit $\unsure{\mathfrak{h}|\mathfrak{b}}$ une sous-algèbre de Cartan [existe tjrs], et considérons $N_G(\mathfrak{h})$, alors l'argument déjà utilisé montre que $N_G(\mathfrak{h}).K=G$, donc on \deleted{peut \unsure{remplacer}} a une suite exacte
% slanted
\margin{\uline{grâce au} \uline{thm de} \uline{conjugaison} \uline{pour ss-algèbres} \uline{de Cartan}}
\diagram{\begin{tikzcd}[ampersand replacement=\&]
e \arrow[r,dash] \& N_K(\mathfrak{h}) \arrow[r] \& N_G(\mathfrak{h}) \arrow[r,dash] \& G' \arrow[r,dash] \& e
\end{tikzcd}}
et on peut donc supposer que $G$ \uline{normalise} $\mathfrak{h}$. Mais alors $K$ normalise $\mathfrak{h}$, donc $\mathfrak{k}$ normalise $\mathfrak{h}$, donc $\mathfrak{h}=\mathfrak{k}$, donc $\mathfrak{k}$ est \uline{nilpotente}. Donc $\mathfrak{k}$ se dévisse \deleted{\unsure{en}} par commutateurs, donc (par la \unsure{Théorie} \unsure{Cartier|Cartan}] $K$ aussi se dévisse en groupes commutatifs [N.B. $[\mathfrak{h},\mathfrak{h}]$ est une sous-$p$-algèbre de Lie \unsure{p-fermée} \unsure{?} correspondant au groupe des \unsure{comm.}]. On peut donc supposer $K$ comm., OK.
% archive sheet number: 8

